
\section{Related Work}
\label{sec:related}

\subsection{Directional Asymmetry in Language Models}
\label{sec:rw-reversal}

Autoregressive models can learn a relation in one direction without learning its reverse: a model trained on ``A is B'' often fails to answer ``B is A''~\cite{berglund2024reversal, grosse2023influence}.
Lv et al. attribute this reversal curse to the next-token objective, which maximizes $p(b \mid a)$ but offers no guarantee on $p(a \mid b)$~\cite{lv2024reversal}, while Wang and Sun trace it to the way transformers represent concepts~\cite{wang2025binding}.
Remedies generally expose the model to both directions during training, either by adding reversed strings to the data~\cite{golovneva2024reverse} or, as in back-translation and its use in TransCoder for code, by generating training pairs from the reverse direction~\cite{sennrich2016backtranslation, lachaux2021unsupervised}.

These studies concern facts or training pairs with a single correct reverse answer.
The inverse tasks we study are executable procedures whose outputs can be tested and admit many correct forms (Definition~\ref{def:pair}), so we score progress toward the source rather than exact recovery.
We also shift the question from pretraining to downstream fine-tuning.
Our balanced bidirectional fine-tuning (Section~\ref{sec:bft}) is close in intent to reverse training, but relies on explicit inverse pairs rather than reversed strings.

\subsection{Side Effects of Fine-Tuning}
\label{sec:rw-forgetting}

Sequential training has been known to overwrite earlier knowledge~\cite{mccloskey1989catastrophic, french1999catastrophic}, and regularization and distillation methods were designed to limit this forgetting~\cite{kirkpatrick2017overcoming, li2018lwf}.
In code models, replay and regularization reduce forgetting under API drift~\cite{weyssow2023continual}, and parameter-efficient methods such as LoRA forget less than full fine-tuning~\cite{biderman2024lora}.
Kotha et al. propose a complementary account in which fine-tuning shifts the task a model infers from its prompt, so that a disguised prompt can restore abilities that appeared lost~\cite{kotha2024implicit}.
Beyond forgetting, fine-tuning can also produce side effects outside the target task, including weaker safety alignment~\cite{qi2024finetuning} and broad misalignment~\cite{betley2025emergent}.

These studies measure tasks unrelated to the fine-tuning objective.
We instead measure its inverse, which shares inputs, outputs, and vocabulary with the trained task, and we use unrelated tasks only as a control for general forgetting.% (Section~\ref{sec:forgetting}).
The account of Kotha et al. is relevant here, since a fine-tuned model may interpret a reverse prompt as a forward request.
We do not evaluate replay or regularization. They remain candidate mitigations outside this study.

\subsection{Code LLMs on Paired Tasks}
\label{sec:rw-paired}

LLMs are now applied to each side of the pairs we study.
They translate code, although their translations can introduce bugs~\cite{pan2024lost, yang2024unitrans}, and they are used to decompile binaries~\cite{tan2024llm4decompile, hu2024degpt} and to obfuscate code~\cite{mohseni2025obfuscate}.
Each of these studies evaluates the direction it targets.
Other work exploits both directions for evaluation: round-trip correctness runs a model forward and backward to check its outputs without labels~\cite{allamanis2024roundtrip}, and CRUXEval asks models to predict outputs from inputs as well as inputs from outputs~\cite{gu2024cruxeval}.
Semantics-preserving transformations are also used to probe the robustness of code models~\cite{yefet2020adversarial, rabin2021generalizability} and to build positive pairs for contrastive learning~\cite{jain2021contracode}.

These works either evaluate a fixed model or ask whether a prediction remains stable under a transformation.
Our question is whether a model can still undo a transformation after it has learned to apply it.
BRF addresses this question by comparing both directions before and after fine-tuning and by reporting the change as a regression.
Our contrastive fine-tuning reuses the positive-pair idea of ContraCode, but only as a probe (Section~\ref{contrastive}) and not as a training recipe.

\subsection{Faithfulness of Reasoning Traces}
\label{sec:rw-faithfulness}

A generated explanation need not reflect the computation behind an answer~\cite{jacovi2020faithfulness}, and chain-of-thought text can misrepresent the actual reason for a prediction~\cite{turpin2023unfaithful, lanham2023measuring}.
We therefore treat reasoning traces as behavioral evidence only (Section~\ref{sec:cs_mechanism}) and make no claim about internal representations.

\medskip
In summary, prior work has studied directional asymmetry on facts, forgetting on unrelated tasks, and the bidirectional use of models that were not adapted to one direction.
Our study connects these lines for paired code transformations, within the scope of one controlled testbed and two additional settings, and with behavioral evidence on mechanism.